% page 50 of the facsimile (image p-049 of the scan)
% block 160.0 x 242.0 mm ; left margin 8.0 mm
\pgc{-12.700mm}{0.000mm}{
\l{\cel{11}42}
\l{}
\l{}
\l{\sou{arteriotomio}.\cel{2}Aperto di arterio per lanceto. - DEFIRS.}
\l{}
\l{\sou{artezo-puteo}.\cel{2}Puteo borita til aquo-strato subsula qua lore}
\l{\cel{3}spriceskas. - DEFIRS.}
\l{}
\l{\sou{artichoko}. (\sou{bot.)}\cel{2}Sorto di kardono di qua la floro konsistas}
\l{\cel{3}ye folii imbrikita, kun fundamento pulpoza. - L.\cel{1}\sou{cynara}}
\l{\cel{3}\sou{scolymus}.\cel{2}- DEFR.}
\l{}
\l{\sou{artifico}.\cel{2}To quo uzesas por travestiar naturalajo, verajo,}
\l{\cel{3}por trompar. - DEFIS.}
\l{}
\l{\sou{artiko}. (\sou{anat.})\cel{2}Junturo naturala inter du parti di korpo,}
\l{\cel{3}movebla ica sur ita. - DEFIS.}
\l{}
\l{\sou{artiklo}. I. (\sou{gram}.)Partikulo qua, en ula lingui nacionala,}
\l{\cel{3}juntesas a substantivo quan lu determinas. - II. (\sou{jurnalo},}
\l{\cel{3}\sou{revuo, e c.}) Singla de la parti di edituro kolektiva.}
\l{\cel{3}- DEFIRS.}
\l{}
\l{\sou{artikular}.\cel{2}(\sou{trans}.) Emisar son-uni vocala, donante a li formi}
\l{\cel{3}diversa, per u\sou{la}\cel{1}movi da la labii, dá la lango.- DEFIRS.}
\l{}
\l{\sou{artilrio}. (\sou{milit-arto)}\cel{2}Armeo-korpo uzata por la funcionigo}
\l{\cel{3}di la kanoni. - DEFIRS.}
\l{}
\l{\sou{artro}. (\sou{anat.})\cel{2}Korpo-parto movebla sur altra a qua lu esas}
\l{\cel{3}juntita. - FS.}
\l{}
\l{\sou{artrito}. (\sou{patol}.)Inflamuro di la tisui fibroza e seroza di}
\l{\cel{3}la artiki. - DEFIRS.}
\l{}
\l{\sou{artrodio}.\cel{2}I. (\sou{anat.)}\cel{2}Artiko en qua esas inkastrita la extrem-}
\l{\cel{3}ajo di altra osto. - II. (\sou{mekan.})\cel{2}Junturo qua su movas}
\l{\cel{3}quale genuo. -\cel{2}EFS.}
\l{}
\l{\sou{artropodo}. (\sou{zool.}) Brancho di la animalo-regno, qua inkluzas}
\l{\cel{3}la animali chitinoza (krabi, kankri, e c.).- DEFIS.}
\l{}
\l{\sou{arufar}.\cel{2}(\sou{trans.})\cel{2}Manuagar kontre la sinso normala la hari}
\l{\cel{3}sur la kapo. - I.}
\l{}
\l{\sou{arumo}. (\sou{bot}.) Planto herbatra di qua ula speci kontenas nutro-}
\l{\cel{3}fekulo. - L.\cel{1}\sou{arum}.\cel{2}- DEFIRS.}
\l{}
\l{\sou{aruspico}. (\sou{che la Romani antiqua}). Sacerdoto qua exploris la}
\l{\cel{3}internaji di la viktimi por obtenar de ta internaji auguri.}
\l{\cel{3}- DEFIRS.}
\l{}
\l{\sou{- as}. (\sou{gram.}) Finalo di la indikativo prezenta (e sen-epoko) di}
\l{\cel{3}la verbi.}
\l{}
\l{\sou{aso}.\cel{2}Ludo-karto o ludo-kubo di nur un punto. - DEFIS.}
}
